\documentclass[UTF8,11pt,a4paper]{ctexart}
\usepackage[a4paper,top=21mm,bottom=23mm,left=24mm,right=24mm]{geometry}
\usepackage{amsmath,amssymb,mathtools,bm}
\usepackage{booktabs,tabularx,array,longtable}
\usepackage{enumitem}
\usepackage{xcolor}
\usepackage{hyperref}
\usepackage{fancyhdr}
\usepackage{setspace}
\usepackage{microtype}
\usepackage[most]{tcolorbox}

\definecolor{MainBlue}{RGB}{35,72,106}
\definecolor{SoftBlue}{RGB}{239,246,251}
\definecolor{SoftGray}{RGB}{246,247,249}
\definecolor{DarkGray}{RGB}{65,70,78}
\hypersetup{colorlinks=true,linkcolor=MainBlue,urlcolor=MainBlue}
\pagestyle{fancy}
\fancyhf{}
\fancyhead[L]{\small FIBRE 数学建模}
\fancyhead[R]{\small 从稀疏关系到双向候选排序}
\fancyfoot[C]{\thepage}
\setlength{\headheight}{14pt}
\setstretch{1.14}
\setlist{itemsep=0.25em,topsep=0.35em,leftmargin=2em}
\setlength{\parindent}{2em}
\setlength{\parskip}{0.15em}

\newcommand{\Rset}{\mathcal R}
\newcommand{\Eset}{\mathcal E}
\newcommand{\Pset}{\mathcal P}
\newcommand{\softplus}{\operatorname{softplus}}
\newcommand{\rank}{\operatorname{rank}}
\newcommand{\diag}{\operatorname{diag}}
\newcommand{\normtwo}[1]{\left\lVert #1\right\rVert_2}

\newtcolorbox{keybox}{
  colback=SoftBlue,
  colframe=MainBlue,
  boxrule=0.55pt,
  arc=1mm,
  left=2.3mm,right=2.3mm,top=1.8mm,bottom=1.8mm,
  before skip=7pt,after skip=7pt,
  fontupper=\small\color{DarkGray}
}

\newtcolorbox{codebox}{
  colback=SoftGray,
  colframe=gray,
  boxrule=0.4pt,
  arc=0.7mm,
  left=2.3mm,right=2.3mm,top=1.8mm,bottom=1.8mm,
  before skip=7pt,after skip=7pt,
  fontupper=\small\ttfamily,
  breakable
}

\title{\vspace{-0.9cm}\Huge\bfseries FIBRE：从稀疏已知关系到双向候选排序\\[0.42em]
\Large 酶--反应检索的数学建模、训练目标与验证}
\author{南京大学二〇二六年国际遗传工程机器大赛团队}
\date{二〇二六年九月二十八日}

\begin{document}
\maketitle

\begin{abstract}
本章只解释当前冻结 FIBRE 检索模型真正计算的量。我们实际解决的问题是：已知少量酶--反应正关联，在其余大量配对均未标注的条件下，如何根据反应分子特征和蛋白序列特征，对候选酶或候选反应进行排序，并且允许训练中没有出现过的新反应和新蛋白直接进入检索。

为解决这一问题，FIBRE 将反应和蛋白分别编码为可计算的连续表示，用一个通用交互分量和八个低维专家分量产生配对分数，再根据查询方向学习不同的专家权重。训练目标由双向多阳性排序损失、直接针对 Top-3/10/20 的排名约束以及专家稳定性正则共同组成。本文从问题定义、输入变量、模型构造、训练目标、排序输出到冻结结果逐步推导，并在每一步给出与当前代码的一一对应。
\end{abstract}

\tableofcontents

\section{先把问题说清楚：我们最终要得到什么}

\subsection{实验任务的输出：一张候选优先级列表}

在反应找酶任务中，研究者给定一个目标反应 $r$，系统需要从候选蛋白集合 $\Eset$ 中返回前 $K$ 个最值得继续验证的酶。在酶找反应任务中，给定一条蛋白 $e$，系统需要从候选反应集合 $\Rset$ 中返回前 $K$ 个反应。

因此，本章要构造的数学对象只有两个排序分数：
\begin{equation}
S_{R\rightarrow E}(r,e),
\qquad
S_{E\rightarrow R}(r,e).
\label{eq:two-scores}
\end{equation}

固定反应 $r$ 时，按 $S_{R\rightarrow E}(r,e)$ 从大到小排列所有候选蛋白；固定蛋白 $e$ 时，按 $S_{E\rightarrow R}(r,e)$ 从大到小排列所有候选反应。

\begin{keybox}
\textbf{这一点决定了整个建模方向。} 当前 FIBRE 分数没有焦耳、千卡每摩尔、秒的倒数等物理单位，也没有被训练成活性概率。它是一个用于相对排序的无量纲学习分数。文档后面的每个公式都直接服务于式 \eqref{eq:two-scores} 的计算或训练。
\end{keybox}

\subsection{为什么这个排序问题难}

当前通用候选宇宙约包含 $185918$ 条蛋白和 $11081$ 条反应，对应
\begin{equation}
185918\times11081=2\,060\,157\,358
\end{equation}
个潜在配对。关系数据库只记录其中极少一部分已知关联。对绝大多数未记录配对，我们只能知道“没有标签”，不能直接知道“实验确认无活性”。

因此，模型同时面对四个具体困难：

\begin{enumerate}[label=\textbf{D\arabic*.}]
\item \textbf{候选规模大。} 无法为二十亿量级的每一个配对单独学习一个参数。
\item \textbf{未标注不等于负例。} 如果把所有未记录配对都当作强负例，容易把真实但尚未记录的关系压低。
\item \textbf{实验只看前几个候选。} 平均分类误差下降，不保证真实阳性进入 Top-10 或 Top-20。
\item \textbf{一种交互表示很难覆盖所有查询。} 同一套反应特征和蛋白特征中可能同时存在多种可用于区分候选的组合模式；若所有查询只能经过一个低维点积分数，模型只能学习一种统一的配对坐标。
\item \textbf{两个检索方向的查询条件不同。} 反应找酶由反应决定“应该更依赖哪种交互模式”；酶找反应则由蛋白决定这一选择。
\end{enumerate}

后面的模型结构分别对应这五个困难。

\begin{keybox}
\textbf{先看完整计算链：}\\[0.3em]
反应 SMILES $\rightarrow \bm x_R\in\mathbb R^{8270}$ $\rightarrow$ 反应塔 $\rightarrow$ 通用表示、八个专家表示、反应门控；\\
蛋白序列 $\rightarrow$ ESM-C $\rightarrow \bm x_E\in\mathbb R^{1152}$ $\rightarrow$ 蛋白塔 $\rightarrow$ 通用表示、八个专家表示、蛋白门控；\\
两侧对应表示做点积 $\rightarrow$ 九个分量分数 $\rightarrow$ 查询侧门控加权 $\rightarrow S_{R\rightarrow E}$ 或 $S_{E\rightarrow R}$；\\
已知关系 $\rightarrow$ 簇屏蔽 $\rightarrow$ 多阳性损失与 Top-K 损失 $\rightarrow$ 更新网络参数；\\
推理时只保留前半条链：输入查询 $\rightarrow$ 计算所有候选分数 $\rightarrow$ 排序 $\rightarrow$ Top-K 实验短名单。
\end{keybox}

\section{建模输入：代码真正读入了哪些变量}

\subsection{反应输入：把一条反应变成 8270 维向量}

冻结模型使用 \texttt{multiview} 反应特征。对每一条规范化反应，代码从底物和产物中构造七块信息：

\begin{center}
\begin{tabularx}{0.94\linewidth}{>{\raggedright\arraybackslash}p{3.3cm} >{\centering\arraybackslash}p{1.5cm} X}
\toprule
特征块 & 维度 & 在模型中的作用 \\
\midrule
整体反应指纹 & 2048 & 描述整条有向反应的局部结构变化模式 \\
底物 Morgan 指纹 & 2048 & 保留主要底物的分子结构信息 \\
产物 Morgan 指纹 & 2048 & 保留主要产物的分子结构信息 \\
有符号差分 $m_p-m_s$ & 2048 & 显式保留“产物相对底物发生了什么变化” \\
分子描述量 & 11 & 碳、氧、磷和环数及其变化 \\
前体类别独热编码 & 10 & 区分不同萜类前体范围 \\
产物骨架类别独热编码 & 57 & 区分链状、单环、双环、多环及氧化状态组合 \\
\midrule
总计 & 8270 & 冻结运行中的实际反应输入维度 \\
\bottomrule
\end{tabularx}
\end{center}

记整体反应指纹为 $d(r)$，底物和产物 Morgan 指纹分别为 $m_s(r),m_p(r)$，十一维描述量为 $c(r)$，两组类别编码为 $u(r),v(r)$。实际反应向量可以写成
\begin{equation}
\bm x_R(r)=
\big[
 d(r),
 m_s(r),
 m_p(r),
 m_p(r)-m_s(r),
 c(r),
 u(r),
 v(r)
\big]\in\mathbb R^{8270}.
\label{eq:reaction-input}
\end{equation}

式 \eqref{eq:reaction-input} 中真正体现反应方向的是整体反应指纹以及 $m_p-m_s$ 的有符号差分。这里不需要额外引入一个没有进入代码的“反应能量”变量。

\subsection{蛋白输入：ESM-C 600M 的 1152 维平均表示}

冻结模型的蛋白侧输入来自 ESM-C 600M。给定氨基酸序列，预训练蛋白语言模型产生残基表示，当前流水线使用平均池化后的 $1152$ 维向量，并做 $L_2$ 归一化。记为
\begin{equation}
\bm x_E(e)\in\mathbb R^{1152}.
\label{eq:protein-input}
\end{equation}

这意味着蛋白能否进入基础评分，取决于能否从序列得到式 \eqref{eq:protein-input}；训练关系表中的实体身份不会参与评分定义。

\subsection{关系标签：只把确认关联记为正例}

令
\begin{equation}
\Pset\subset\Rset\times\Eset
\end{equation}
表示训练数据中记录的已知酶--反应关联。定义观测矩阵
\begin{equation}
Y_{ij}=
\begin{cases}
1, & (r_i,e_j)\in\Pset,\\
?, & (r_i,e_j)\notin\Pset.
\end{cases}
\label{eq:pu-label}
\end{equation}

问号表示“未标注”。当前模型没有把它改写成实验负例 $0$。这一标签语义直接决定后面的分母屏蔽和多阳性损失。

\section{第一步建模：用因子化双塔解决大候选空间和新实体输入}

\subsection{每一侧先学习自己的连续函数}

反应和蛋白分别经过一个投影塔。两侧结构相同，参数独立。以任意输入 $\bm x$ 为例，共享骨干计算
\begin{equation}
\bm h
=
\operatorname{Dropout}
\left(
\operatorname{GELU}
\left(W_h\operatorname{LN}(\bm x)+\bm b_h\right)
\right),
\label{eq:backbone}
\end{equation}
其中冻结配置的隐藏维度为 $512$，Dropout 为 $0.1$。

随后从 $\bm h$ 产生一个通用表示和八个专家表示：
\begin{align}
\bm z^{(0)} &= \operatorname{norm}(W_0\bm h+\bm b_0)
\in\mathbb R^{128},\\
\bm z^{(k)} &= \operatorname{norm}(W_k\bm h+\bm b_k)
\in\mathbb R^{32},\qquad k=1,\ldots,8,
\label{eq:embeddings}
\end{align}
其中
\begin{equation}
\operatorname{norm}(\bm a)=\frac{\bm a}{\normtwo{\bm a}}.
\end{equation}

对反应和蛋白分别记为
\begin{equation}
\bm z_R^{(k)}(r),\qquad \bm z_E^{(k)}(e).
\end{equation}

\subsection{每个分量只做一个真实存在的计算：点积}

通用分量的配对分数是
\begin{equation}
s_0(r,e)=
\bm z_R^{(0)}(r)^{\mathsf T}
\bm z_E^{(0)}(e),
\label{eq:global-score}
\end{equation}
八个专家分量分别为
\begin{equation}
s_k(r,e)=
\bm z_R^{(k)}(r)^{\mathsf T}
\bm z_E^{(k)}(e),
\qquad k=1,\ldots,8.
\label{eq:expert-score}
\end{equation}

由于所有向量都做了单位范数归一化，九个分量都满足
\begin{equation}
-1\le s_k(r,e)\le1.
\label{eq:component-bound}
\end{equation}

这里的八个专家是从同一组冻结分子输入中学习出的八个低维交互子空间。当前冻结运行没有把它们预先指定成“质子转移专家”“金属专家”等物理机制，也没有给八个专家绑定八种外部数据模态。

\subsection{为什么这种因子化直接解决 D1}

如果为每个潜在配对直接学习一个独立分数，需要面对约 $2.06\times10^9$ 个位置，而且一个新反应或新蛋白没有对应的训练位置。双塔改成学习两个函数
\begin{equation}
F_R:\bm x_R(r)\mapsto
\left(\bm z_R^{(0)},\ldots,\bm z_R^{(8)}\right),
\end{equation}
\begin{equation}
F_E:\bm x_E(e)\mapsto
\left(\bm z_E^{(0)},\ldots,\bm z_E^{(8)}\right).
\end{equation}

候选分数由式 \eqref{eq:global-score}--\eqref{eq:expert-score} 即时计算。模型学习从分子输入到低维向量的函数，由这些函数即时生成配对分数，从而避免为二十亿个潜在配对分别维护独立参数。

\begin{keybox}
\textbf{双未见能力由这里产生。} 一个新反应 $r^*$ 和一个新蛋白 $e^*$ 即使都没有出现在训练关系表，只要能计算 $\bm x_R(r^*)$ 和 $\bm x_E(e^*)$，九个 $s_k(r^*,e^*)$ 就全部有定义。这个结论来自模型计算图本身，不需要额外的流形假设。
\end{keybox}

\section{第二步建模：用方向化专家权重解决“一套交互模式不够用”}

上一节若只保留通用分量 $s_0$，所有查询都只能在同一组 $128$ 维坐标中比较候选。FIBRE 在同一个共享隐藏表示上再学习八组 $32$ 维配对坐标，让不同查询可以强调不同的低维交互模式。这个设计是否真的有用，不能由公式本身宣布；第十节会用“同一训练模型的完整读出 vs. 只保留通用分量”直接检验。

\subsection{模型怎样为不同查询选择专家}

反应塔和蛋白塔各自从隐藏层生成八维门控：
\begin{align}
\bm q^R(r)
&=operatorname{softmax}
\left(\frac{W_g^R\bm h_R(r)+\bm b_g^R}{T_g}\right),\\
\bm q^E(e)
&=operatorname{softmax}
\left(\frac{W_g^E\bm h_E(e)+\bm b_g^E}{T_g}\right),
\label{eq:gates}
\end{align}
冻结配置中 $T_g=1$，并且
\begin{equation}
q_k^R\ge0,quad \sum_{k=1}^{8}q_k^R=1,
\qquad
q_k^E\ge0,quad \sum_{k=1}^{8}q_k^E=1.
\end{equation}

模型还学习两个方向各自的专家总权重：
\begin{equation}
\mu_R=\sigma(a_R),
\qquad
\mu_E=\sigma(a_E),
\label{eq:mix-mass}
\end{equation}
初始值都设为 $0.5$，之后由训练更新。

\subsection{反应找酶：由反应决定专家组合}

给定反应 $r$ 时，最终分数为
\begin{equation}
\boxed{
S_{R\rightarrow E}(r,e)
=
(1-\mu_R)s_0(r,e)
+
\mu_R\sum_{k=1}^{8}q_k^R(r)s_k(r,e)
}
\label{eq:r2e-score}
\end{equation}

这里同一个反应对所有候选蛋白使用同一组 $q^R(r)$。因此排序差异来自各候选在九个交互分量上的得分，而专家选择只由查询反应决定。

\subsection{酶找反应：由蛋白决定专家组合}

给定蛋白 $e$ 时，模型使用蛋白侧门控：
\begin{equation}
\boxed{
S_{E\rightarrow R}(r,e)
=
(1-\mu_E)s_0(r,e)
+
\mu_E\sum_{k=1}^{8}q_k^E(e)s_k(r,e)
}
\label{eq:e2r-score}
\end{equation}

式 \eqref{eq:r2e-score} 和式 \eqref{eq:e2r-score} 共享同一组九个基础配对分数，只改变“当前查询应该怎样组合它们”。这就是代码中的 directional multi-expert。

\subsection{这个结构有三个可以直接验证的数学性质}

第一，两个方向的权重都严格求和为 $1$：
\begin{equation}
(1-\mu_R)+\mu_R\sum_{k=1}^{8}q_k^R=1,
\end{equation}
酶找反应同理。

第二，由式 \eqref{eq:component-bound} 和凸组合可得
\begin{equation}
-1\le S_{R\rightarrow E}(r,e)\le1,
\qquad
-1\le S_{E\rightarrow R}(r,e)\le1.
\label{eq:score-bound}
\end{equation}

第三，整个分数矩阵仍然保持因子化。设通用反应/蛋白嵌入矩阵为 $Z_R^{(0)},Z_E^{(0)}$，第 $k$ 个专家矩阵为 $Z_R^{(k)},Z_E^{(k)}$。反应找酶矩阵可写成
\begin{equation}
S_R
=
(1-\mu_R)Z_R^{(0)}Z_E^{(0)\mathsf T}
+
\mu_R\sum_{k=1}^{8}
\diag(\bm q_k^R)
Z_R^{(k)}Z_E^{(k)\mathsf T}.
\label{eq:r2e-matrix}
\end{equation}
因此
\begin{equation}
\operatorname{rank}(S_R)
\le128+8\times32=384.
\label{eq:rank-bound}
\end{equation}
酶找反应也有同样的上界，只是查询权重作用在蛋白一侧。

\begin{keybox}
式 \eqref{eq:rank-bound} 是对当前网络结构的直接代数结论。它说明模型可以表达多组交互模式，同时仍通过低维因子计算大规模候选分数；这里不需要假设存在某个不可观测的物理“真实核”。
\end{keybox}

\section{第三步建模：让未标注配对保持“未知”}

\subsection{为什么训练分母需要屏蔽}

对反应 $r_i$，记它的已知阳性蛋白集合为
\begin{equation}
P_R(i)=\{j:(r_i,e_j)\in\Pset\}.
\end{equation}
对蛋白 $e_j$，记其已知阳性反应集合为
\begin{equation}
P_E(j)=\{i:(r_i,e_j)\in\Pset\}.
\end{equation}

直接把其余所有配对都当作训练负例会产生明显问题：高度相似的蛋白可能催化相同反应，高度相似的反应也可能由同一蛋白催化，但数据库未必把所有这些关系都记录齐全。

冻结训练利用预先构造的蛋白簇和反应簇做保守屏蔽。对反应找酶，如果一个未标注蛋白与某个已知阳性蛋白属于同一蛋白簇，就不让这个未标注配对进入该查询的对比学习分母。记允许的分母集合为 $D_R(i)$。酶找反应同理得到 $D_E(j)$。

数学上只需要满足
\begin{equation}
P_R(i)\subseteq D_R(i),
\qquad
P_E(j)\subseteq D_E(j),
\end{equation}
并且部分高风险未标注项被排除。

\begin{keybox}
这一步没有声称被屏蔽的配对一定是阳性，也没有把未屏蔽配对解释成实验确认负例。它只是改变哪些未标注项有资格在训练时充当对比项。
\end{keybox}

\section{第四步建模：直接优化双向、多阳性的检索顺序}

\subsection{双向多阳性损失}

将反应找酶分数除以温度 $\tau=0.07$：
\begin{equation}
\ell^R_{ij}
=
\frac{S_{R\rightarrow E}(r_i,e_j)}{\tau}.
\end{equation}

对一个反应查询，训练损失为
\begin{equation}
L_R(i)
=
\log\sum_{j\in D_R(i)}e^{\ell^R_{ij}}
-
\log\sum_{j\in P_R(i)}e^{\ell^R_{ij}}.
\label{eq:r2e-contrastive}
\end{equation}

第一项把所有允许竞争的候选放入分母，第二项把该查询的所有已知阳性共同放入分子。一个反应有多个已知酶时，不需要人为挑选唯一正确答案。

酶找反应方向完全对应：
\begin{equation}
L_E(j)
=
\log\sum_{i\in D_E(j)}e^{\ell^E_{ij}}
-
\log\sum_{i\in P_E(j)}e^{\ell^E_{ij}},
\label{eq:e2r-contrastive}
\end{equation}
其中
\begin{equation}
\ell^E_{ij}=\frac{S_{E\rightarrow R}(r_i,e_j)}{\tau}.
\end{equation}

冻结配置对两个方向等权：
\begin{equation}
L_{\mathrm{ctr}}
=
\frac12L_R+\frac12L_E.
\label{eq:contrastive-total}
\end{equation}

\subsection{为什么还要再加一个 Top-K 损失}

式 \eqref{eq:contrastive-total} 提高阳性相对分数，但实验真正关心的是“至少有一个阳性能否进入前 $K$”。因此模型额外直接比较最佳阳性和第 $K$ 个允许负例。

对任意查询 $q$，令
\begin{equation}
p(q)=\max_{a\in P(q)}\ell(q,a)
\end{equation}
为最佳阳性 logit，令
\begin{equation}
n_K(q)
\end{equation}
为允许参与对比的未标注候选中第 $K$ 大的 logit。若 $n_K(q)>p(q)$，说明至少有 $K$ 个当前竞争候选排在最佳阳性前面。这里的“竞争候选”仍然保持未标注语义，不被解释成实验负例。

对应平滑惩罚为
\begin{equation}
L_K(q)=\softplus\big(n_K(q)-p(q)\big).
\label{eq:topk-one}
\end{equation}

冻结模型同时关注 $K=3,10,20$：
\begin{equation}
L_{\mathrm{TopK}}
=
\frac{
0.10L_3+0.05L_{10}+0.025L_{20}
}{0.175},
\label{eq:topk-combine}
\end{equation}
再对两个检索方向等权平均。

\begin{keybox}
这是模型和实验预算之间最直接的数学连接。Top-3、Top-10、Top-20 在这里直接进入训练目标，因此模型优化的对象就是研究者真正会取出的短名单位置。
\end{keybox}

\section{第五步建模：让八个专家形成分工，同时避免失控}

如果只有式 \eqref{eq:r2e-score}--\eqref{eq:e2r-score}，门控可能长期只用少数专家，不同专家也可能学习出几乎相同的表示。冻结训练加入四类正则。

\subsection{批次层面的专家使用平衡}

设一个 batch 内门控均值为
\begin{equation}
\bar q_k=\frac1N\sum_{n=1}^{N}q_{nk}.
\end{equation}
八个专家的平衡损失为
\begin{equation}
L_{\mathrm{bal}}
=
8\sum_{k=1}^{8}\left(\bar q_k-\frac18\right)^2.
\label{eq:balance}
\end{equation}
反应侧和蛋白侧分别计算后取平均。它抑制所有查询长期集中到同一个专家。

\subsection{单查询门控尽量明确}

定义平均门控熵
\begin{equation}
H_q
=
-\frac1N\sum_{n=1}^{N}
\sum_{k=1}^{8}q_{nk}\log q_{nk}.
\label{eq:entropy}
\end{equation}

当前总损失对 $H_q$ 使用正权重，因此最小化训练会鼓励单个查询的门控更集中。它和式 \eqref{eq:balance} 的作用层次不同：平衡项要求整个 batch 不要只用一个专家，熵项要求每个具体查询尽量做出明确选择。

\subsection{专家表示减少重复}

对同一个对象的八个单位专家向量，构造 Gram 矩阵
\begin{equation}
G_{k\ell}
=
\bm z^{(k)\mathsf T}\bm z^{(\ell)}.
\end{equation}
多样性惩罚为所有非对角项平方的平均：
\begin{equation}
L_{\mathrm{div}}
=
\operatorname{mean}_{k\ne\ell}G_{k\ell}^2.
\label{eq:diversity}
\end{equation}
它鼓励不同专家使用不同的低维方向，减少八个分量完全复制同一个表示。

\subsection{高权重分量之间不应任意矛盾}

把通用分量和八个专家统一编号为 $a=0,\ldots,8$。对某个方向，记它们的组合权重为 $\rho_a$，对应分数为 $s_a$。当前实现对同时具有正权重的分量加入
\begin{equation}
C(\bm s,\bm\rho)
=
\frac{
\sum_{a<b}\rho_a\rho_b(s_a-s_b)^2
}{
\sum_{a<b}\rho_a\rho_b
},
\label{eq:consistency}
\end{equation}
再对所有训练配对求平均，并对两个方向等权组合成 $L_{\mathrm{con}}$。

式 \eqref{eq:consistency} 的含义很直接：如果两个分量都被门控赋予较高权重，它们对同一配对给出的分数差异过大就会受到惩罚。冻结配置的权重很小，只取 $0.02$，因此它承担弱稳定约束。

\section{完整模型：所有真正生效的项放在一张式子里}

冻结 FIBRE 复现实验中，机制辅助头和机制额外分数权重均为 $0$，hard-negative 子采样也关闭。因此当前真正优化的目标可以完整写成
\begin{equation}
\boxed{
L
=
L_{\mathrm{ctr}}
+
L_{\mathrm{TopK}}
+
0.05L_{\mathrm{bal}}
+
0.005H_q
+
0.01L_{\mathrm{div}}
+
0.02L_{\mathrm{con}}
}
\label{eq:full-loss}
\end{equation}

这里 $L_{\mathrm{bal}},H_q,L_{\mathrm{div}}$ 都由反应塔和蛋白塔对应项平均得到。

训练采用 AdamW，学习率 $3\times10^{-4}$，权重衰减 $10^{-4}$，共训练 $80$ 个 epoch，梯度范数裁剪到 $5.0$，随机种子固定为 $20260723$。

\begin{center}
\begin{tabularx}{0.92\linewidth}{>{\raggedright\arraybackslash}p{4.4cm} >{\raggedright\arraybackslash}p{3.0cm} X}
\toprule
参数 & 冻结值 & 作用 \\
\midrule
隐藏层维度 & 512 & 两侧共享骨干宽度 \\
通用交互维度 & 128 & 基础配对分量 \\
专家数量 & 8 & 查询条件化的交互子空间数量 \\
单专家维度 & 32 & 每个专家的低维配对宽度 \\
门控温度 & 1.0 & softmax 门控尺度 \\
专家混合初值 & 0.5 & $\mu_R,\mu_E$ 的 sigmoid 初值 \\
对比温度 & 0.07 & 排序 logit 尺度 \\
方向权重 & 0.5 / 0.5 & R2E 与 E2R 等权 \\
Top-K 权重 & 0.10 / 0.05 / 0.025 & 分别对应 3 / 10 / 20 \\
平衡 / 熵 / 多样性 & 0.05 / 0.005 / 0.01 & 控制专家使用与去冗余 \\
一致性权重 & 0.02 & 约束高权重分量的分数冲突 \\
\bottomrule
\end{tabularx}
\end{center}

\section{从分数到实验短名单：推理过程只有六步}

对一个反应找酶查询，运行过程可以压缩为：

\begin{enumerate}[label=\textbf{Step \arabic*.}]
\item 从反应 SMILES 构造式 \eqref{eq:reaction-input} 的 $8270$ 维向量。
\item 通过反应塔得到 $1$ 个通用向量、$8$ 个专家向量和反应门控 $\bm q^R(r)$。
\item 候选蛋白从预计算的 ESM-C 表示进入蛋白塔，得到对应的九组向量。
\item 用式 \eqref{eq:global-score} 和 \eqref{eq:expert-score} 计算九个分量分数。
\item 用式 \eqref{eq:r2e-score} 合成为每个候选的最终 R2E 分数。
\item 按分数降序排列，返回前 $K$ 个候选。
\end{enumerate}

酶找反应完全相同，只把查询侧换成蛋白并使用式 \eqref{eq:e2r-score}。

\begin{codebox}
代码对应：\\
\texttt{ExpertTower.forward}: 式 (\ref{eq:backbone})--(\ref{eq:gates})\\
\texttt{DirectionalMultiExpertDualTower.score\_matrices}: 式 (\ref{eq:global-score})--(\ref{eq:e2r-score})\\
\texttt{directional\_multi\_positive\_loss}: 式 (\ref{eq:r2e-contrastive})--(\ref{eq:contrastive-total})\\
\texttt{directional\_topk\_surrogate}: 式 (\ref{eq:topk-one})--(\ref{eq:topk-combine})\\
\texttt{gate\_regularization}: 式 (\ref{eq:balance})--(\ref{eq:diversity})\\
\texttt{overlap\_consistency\_loss}: 式 (\ref{eq:consistency})
\end{codebox}

\section{如何判断模型有没有解决问题}

\subsection{评测指标直接围绕候选顺序}

对一个查询 $q$，如果有多个已知阳性，定义最佳阳性名次
\begin{equation}
r^*(q)=\min_{a\in P(q)}\rank_q(a).
\label{eq:best-rank}
\end{equation}

平均倒数排名为
\begin{equation}
\mathrm{MRR}
=
\frac1N\sum_{q=1}^{N}\frac1{r^*(q)}.
\label{eq:mrr}
\end{equation}

Top-$K$ 命中概率为
\begin{equation}
\mathrm{Hit@K}
=
\frac1N\sum_{q=1}^{N}
\mathbf 1\big[r^*(q)\le K\big].
\label{eq:hitk}
\end{equation}

当一个查询有多个阳性时，还计算
\begin{equation}
\mathrm{PositiveRecall@K}(q)
=
\frac{|\operatorname{TopK}(q)\cap P(q)|}{|P(q)|}.
\end{equation}

这些指标和训练中的式 \eqref{eq:topk-one} 使用相同的排序对象。

\subsection{协议一：熟悉候选范围下，模型能否给出有用的短名单}

在 TPS \texttt{legacy\_exact} 协议中共有 $513$ 个反应查询。冻结模型得到：

\begin{center}
\begin{tabular}{lrrrr}
\toprule
协议 & MRR & Hit@10 & Hit@20 & 最佳阳性中位名次 \\
\midrule
TPS legacy exact, R2E & 0.23695 & 35.87\% & 45.81\% & 29 \\
\bottomrule
\end{tabular}
\end{center}

这个协议回答的是：在熟悉的 TPS 候选范围中，式 \eqref{eq:r2e-score} 是否能把已知阳性放进较短实验列表。

\subsection{协议二：反应和蛋白同时冷启动时，函数因子化是否还能工作}

严格 \texttt{double\_cold\_25cell} 把反应和蛋白按独立折划分，训练时同时排除测试反应折和测试蛋白折对应的关系；评测覆盖所有严格折组合。冻结结果为：

\begin{center}
\begin{tabular}{lrrrr}
\toprule
方向 & MRR & Hit@10 & Hit@20 & 最佳阳性中位名次 \\
\midrule
R2E & 0.04489 & 10.78\% & 19.05\% & 136 \\
E2R & 0.07229 & 19.02\% & 28.71\% & 64 \\
\bottomrule
\end{tabular}
\end{center}

这组结果直接检验第三节所说的新实体输入能力：测试查询和对应候选关系不能靠训练 ID 查表获得，必须通过分子特征函数产生分数。

\subsection{协议三：扩大到 185,918 条蛋白后，八个专家是否真的贡献排序信息}

最直接的结构消融是在同一个已训练模型上比较完整多专家读出和只使用通用分量。513 个 TPS 反应的严格双冷启动查询被放入完整的 $185918$ 蛋白候选宇宙中：

\begin{center}
\begin{tabular}{lrrrr}
\toprule
模型读出 & MRR & Hit@10 & Hit@20 & 最佳阳性中位名次 \\
\midrule
仅通用分量 & 0.009838 & 2.521\% & 4.622\% & 8998 \\
完整方向化多专家 & 0.012448 & 3.221\% & 5.602\% & 5488 \\
\bottomrule
\end{tabular}
\end{center}

完整模型相对通用分量的 MRR 增加 $0.002610$，相对提升约 $26.5\%$；Hit@10 增加约 $0.70$ 个百分点，Hit@20 增加约 $0.98$ 个百分点，最佳阳性中位名次前移 $3510$ 位。

这项消融回答第五节之前最关键的问题：八个专家和方向化门控是否只是增加参数。相同训练模型、相同查询和相同 $185918$ 候选下，完整读出优于只保留通用分量，说明多专家分解确实改变了大候选空间的排序。

\subsection{结果也显示模型仍有明显困难}

广域候选中的绝对 Hit@10 仍只有 $3.22\%$，严格双冷启动也显著低于熟悉协议。模型解决了“任何有效输入可以进入评分”和“专家分解能在大候选范围提供额外排序信息”两个问题，没有把开放世界酶功能预测变成一个已经饱和的任务。

历史八专家路线与当前冻结版本的比较也存在方向性取舍：当前版本在 TPS R2E 的 MRR 和 Hit@20 上提高，Hit@10 约下降 $1.45$ 个百分点；E2R 的 MRR、Hit@10 和 Hit@20 则明显提高。这也是保留双方向独立评测的原因。

\section{把整个建模逻辑重新连成一条线}

现在可以从实验问题一路走到最终公式，中间没有额外的物理故事：

\begin{enumerate}[label=\textbf{\arabic*.}]
\item \textbf{实验只能测少量候选。} 因此输出定义成排序，核心指标是最佳阳性名次和 Hit@K。
\item \textbf{候选对太多。} 因此不用配对 ID 矩阵，而从反应分子特征和蛋白序列特征分别学习低维函数表示，再做点积。
\item \textbf{新反应和新蛋白也要能进来。} 因此所有分数都由原始分子输入经过函数计算，训练 ID 不进入评分公式。
\item \textbf{一组低维交互不足以覆盖所有查询。} 因此从共享隐藏表示中学习一个通用分量和八个专家分量，并让查询侧门控决定组合权重。
\item \textbf{两个检索方向的“查询”不同。} 因此 R2E 使用反应门控，E2R 使用蛋白门控，得到两个方向化分数。
\item \textbf{未标注关系不能粗暴当成实验负例。} 因此利用蛋白簇和反应簇把部分高风险未标注项从训练分母中屏蔽。
\item \textbf{实验真正关心前 3、10、20。} 因此除多阳性对比损失外，再直接惩罚最佳阳性落在第 $K$ 个竞争候选之后的情况。
\item \textbf{多专家可能塌缩或重复。} 因此加入批次平衡、单查询低熵、专家去冗余和弱一致性约束。
\item \textbf{最终用真实候选排序验证。} 熟悉协议、严格双冷启动和 $185918$ 蛋白广域消融分别检验短名单能力、新实体泛化和多专家的实际贡献。
\end{enumerate}

\begin{keybox}
这就是当前 FIBRE 数学模型的完整闭环：\textbf{分子输入 $\rightarrow$ 因子化表示 $\rightarrow$ 多专家配对 $\rightarrow$ 方向化组合 $\rightarrow$ PU-aware 排序训练 $\rightarrow$ Top-K 短名单 $\rightarrow$ 严格协议验证。} 每个箭头都在当前代码和冻结实验中真实存在。
\end{keybox}

\section{建模边界：当前公式能说明什么}

当前模型可以支持以下陈述：

\begin{itemize}
\item FIBRE 学习的是酶--反应候选的相对排序分数；
\item 分数由反应分子特征和蛋白序列表示直接计算，因此有效新输入具有评分定义；
\item 当前冻结模型使用一个 $128$ 维通用分量和八个 $32$ 维专家分量；
\item 两个方向共享分量分数，查询侧门控和专家总权重独立；
\item 训练显式使用多阳性、簇屏蔽和 Top-3/10/20 目标；
\item 广域同模型消融显示专家组合相对通用分量产生可测排序增益。
\end{itemize}

当前模型没有直接支持以下解释：

\begin{itemize}
\item 排序分数的数值差可以直接解释为实验效应大小；
\item 八个专家分别对应八个已命名催化机制；
\item 未记录配对就是负例；
\item 双未见“可计算”等于双未见“一定准确”。
\end{itemize}

这些边界让数学说明和真正运行的模型保持一致。

\section{代码与证据入口}

\begin{longtable}{>{\raggedright\arraybackslash}p{4.0cm} >{\raggedright\arraybackslash}p{10.0cm}}
\toprule
内容 & 仓库入口 \\
\midrule
\endfirsthead
\toprule
内容 & 仓库入口 \\
\midrule
\endhead
反应多视图特征与基础双塔 & \path{projects/active/fibre/runtime/base_model.py} \\
冻结方向化多专家模型 & \path{reproducibility/bime_rank/support/evaluate_multi_expert_protocol_comparison.py} \\
簇屏蔽、多阳性损失、Top-K 损失 & 同上 \\
一致性损失实现 & \path{projects/active/fibre/kernel/atlas.py} 中 \texttt{overlap\_consistency\_loss} \\
冻结参数 & \path{reproducibility/bime_rank/configs/fibre_interaction_atlas_v1.yaml} \\
冻结运行摘要 & \path{results/fibre_interaction_atlas_full_v1/summary.json} \\
综合结果与广域消融 & \path{reproducibility/bime_rank/records/FIBRE_INTERACTION_ATLAS_SCORECARD_V1.json} \\
广域 $185918$ 蛋白评测 & \path{results/fibre_interaction_atlas_tps_broad_universe_v1/} \\
\bottomrule
\end{longtable}

\section{结论}

FIBRE 的数学任务可以用一句话概括：
\begin{keybox}
\textbf{从稀疏已知正关系中学习两个由分子输入定义的方向化排序函数，使真实阳性尽可能进入有限实验预算对应的 Top-K。}
\end{keybox}

模型为此做了五件事：用双塔因子化避免学习完整配对矩阵；用一个通用分量和八个专家分量提高交互表达能力；用方向化门控分别处理反应找酶和酶找反应；用簇屏蔽降低未标注关系被误伤的风险；用多阳性损失和 Top-K 约束把训练目标对准真实检索任务。

这套建模把项目真正面对的五个困难逐一变成了可以计算、训练和验证的数学结构。整份主建模只保留当前计算链中真实使用的变量、公式和约束。

\end{document}
